\documentclass[conference]{IEEEtran}

\usepackage[T1]{fontenc}
\usepackage[utf8]{inputenc}
\usepackage[brazil]{babel}
\usepackage{cite}
\usepackage{microtype}

\title{Pipeline de Execução, Ponto Flutuante e Superescalar}

\author{\IEEEauthorblockN{Autor}
\IEEEauthorblockA{Universidade de Brasília\\
}}

\begin{document}

\maketitle

\section{AMD Instinct MI300}

A série AMD Instinct MI300 emprega a microarquitetura CDNA~3 e organiza o
processamento gráfico em \emph{Accelerator Complex Dies} (XCDs). Cada XCD
possui 40 \emph{Compute Units} (CUs). Quatro \emph{Asynchronous Compute Engines} distribuem grupos de trabalho para essas
CUs. Cada CU constitui um núcleo paralelo completo, com busca e escalonamento de
instruções, unidades escalares, vetoriais e matriciais e unidades de
load/storage \cite{amd-cdna3-whitepaper}. A unidade vetorial (VALU) contém quatro
processadores SIMD de 16 caminhos e arquivos de registradores vetoriais \cite{amd-pipelines}. 
Desvios são tratados por uma máscara de execução, ou seja, executando os dois caminhos possíveis de cada desvio. 
As vias que não pertencem ao caminho tomado tornam-se inativas, reduzindo a
utilização em caso de divergência. A unidade escalar (SALU), é usada para
aritmética, comparações e controle uniformes a toda a onda. As unidades
VMEM e SMEM são usadas para acessos vetoriais e escalares à memória. Também há: a Local Data Share (LDS), para
comunicação local, a unidade de desvios, para executar jumps e branches, e os aceleradores MFMA (\emph{Matrix Fused Multiply--Add}) \cite{amd-cu,amd-pipelines}.
Emr relação às unidades de ponto flutuante, o MI300 é capaz de realizar operações de ponto flutuante de baixa precisão com FP8, FP16 e BF16, além das operações convencionais de 32 bits.

\bibliographystyle{IEEEtran}
\bibliography{referencias}

\end{document}
